\section{Related Research}\label{sec:related-research}

This section reviews the three bodies of work most directly related to the
Intent Recognition Engine---language-model copilots for security operations,
attacks on applications that integrate language models, and methods for
detecting model fabrication---followed by a brief assessment of what each body
contributes to and leaves open for this dissertation.

\subsection*{Language-model copilots for security operations}

Two systems define the current shape of LLM-assisted offensive security. Happe
and Cito proposed language models as \emph{sparring partners} for penetration
testing and demonstrated a closed-loop system in which a model drove a
deliberately vulnerable virtual machine to root access \parencite{happe2023pwned}.
PENTESTGPT treats penetration testing as a task for a language-model-centric
assistant, decomposing an engagement across a reasoning chain, a generation
chain and a guidance chain so that each phase of an assessment is planned and
executed incrementally \parencite{deng2024pentestgpt}. Both contributions
establish that language models can produce \emph{plausible} security actions
and can operate tooling in controlled environments, and both are representative
of the generation-side focus of the field: output is produced for a human
operator to review, and correctness is handled informally. Neither system
subjects its output to deterministic schema, scope and authorisation validation
before machine consumption, which is precisely the gap that a copilot feeding
an \emph{automated} toolchain must close. The IRE inverts this priority: it
treats the model's generation as a candidate that must pass checkable
predicates (Section~\ref{sec:domain-foundations}) and only then becomes an
actionable intent.

\subsection*{Attacks on language-model-integrated applications}

The second body of work establishes that unvalidated model output is dangerous
in integrated systems. Liu \emph{et al.} formalised prompt-injection attacks
against LLM-integrated applications and systematically benchmarked a set of
common defences, finding that prevention-based defences were weak and
detection-based defences incomplete \parencite{liu2023promptinjection}.
Greshake \emph{et al.} demonstrated \emph{indirect} prompt injection, in which
an attacker plants instructions inside content---including web pages and
documents---that an integration later retrieves and feeds to the model, and
showed compromise of real-world assistant-style integrations
\parencite{greshake2023indirect}. Kang \emph{et al.} showed that dual-use
attacks that combine well-known computer-security techniques bypass in-the-wild
content safeguards of deployed models with near-complete success
\parencite{kang2024dualuse}. The OWASP Top~10 for LLM applications
consolidates these classes---prompt injection first among them---as
first-class application risks rather than theoretical concerns
\parencite{owasp2025}.

The implication for this dissertation is direct. A security copilot that
consumes an utterance, asks a model to classify it, and dispatches the result
is exactly the kind of integration these attacks exploit: a single adversarial
utterance can steer the classification output, and any downstream consumer will
then act on the attacker-chosen intent unless an independent, non-model check
intervenes. The relevant design response therefore places deterministic
authorisation and validation \emph{after} generation, so that even a perfectly
diverted model cannot manufacture an accepted intent that violates syntax,
scope or role policy (the safety invariant introduced in
Chapter~\ref{ch:product-design}).

\subsection*{Detecting fabrication and measuring uncertainty}

The third body of work supplies the theoretical toolkit for catching fabricated
output. Ji \emph{et al.} provide the taxonomy of hallucination in natural
language generation that distinguishes intrinsic from extrinsic fabrication
\parencite{ji2023hallucination} (Section~\ref{sec:domain-foundations}).
Sampling-based detection uses the model's own output statistics:
SelfCheckGPT scores the consistency of sampled responses to flag ungrounded
claims \parencite{manakul2023selfcheckgpt}; semantic-uncertainty methods embed
responses and cluster meaning-equivalent outputs to measure disagreement
\parencite{kuhn2023semantic}, and semantic entropy applies this principle to
detect confabulations in free-form text \parencite{farquhar2024semantic}.
Self-consistency decoding accumulates agreement across samples and rewards the
majority answer \parencite{wang2023selfconsistency}.

What these methods share is an implicit budget: they assume either interactive
generation with a user, or the freedom to produce several responses per query
and cluster them. Each additional sample multiplies the inference cost of a
CPU-only model on the operator's laptop, so the design of this dissertation
must transpose the disagreement intuition into a \emph{single} guarded
classification: restricted sampling (a bounded number of responses), a cheap
embedding-based agreement gate, and a deterministic escalation path that does
not depend on a large verifier model.

\subsection*{Grounding via retrieval and embeddings}

A complementary line of work grounds model answers in external corpora rather
than in self-consistency. Retrieval-augmented generation (RAG) conditions
generation on retrieved passages so that answers can be traced to sources
\parencite{lewis2020rag}, and sentence embeddings provide the similarity
signals such systems rely on \parencite{reimers2019sbert}. RAG is attractive
for factual Q\&A, but it introduces retrieved content into the prompt at
inference time---an additional prompt-injection surface
\parencite{greshake2023indirect}---and presupposes corpus infrastructure and
embedding services that contradict the offline, single-device boundary adopted
here (Section~\ref{sec:scope-limitations-assumptions}). The IRE therefore
performs no retrieval at inference time and uses embeddings only for the
agreement and cache signals.

\subsection*{Assessment}

Taken together, the related research establishes three things: copilots can
generate plausible security actions but do not validate them for machine
consumption; integrated systems are exploitable through adversarial and
indirect-injection utterances, and published defences are brittle; and
fabrication can be detected in principle through sampling and semantic
agreement, at a cost that grows with the number of samples. What the surveyed
literature does not measure, to the best of the review conducted here, is the
behaviour of a \emph{guardrail stack}---a full deterministic validation
pipeline with authorisation checks and a bounded self-consistency gate over a
small offline model---under a strict single-inference latency budget, nor the
isolated contribution of each guard layer, nor whether the achieved safety
transfers between different offline model backends. Those measurements are
precisely RQ1--RQ3, and Section~\ref{sec:gap-synthesis} consolidates the gap
they define into the conceptual framework used in the remainder of this
dissertation.

\TODO{At the final AUDIT step, re-verify the characterisation of
\citeauthor{liu2023promptinjection} (defence benchmark) and
\citeauthor{kang2024dualuse} (dual-use bypass) against the full text of their
respective publications, since the descriptions here are relayed through the
project's research paper rather than a direct reading of each source.}